% !TEX root = ../../main.tex
\chapter*{线性代数概念复查表}
\phantomsection
\addcontentsline{toc}{chapter}{线性代数概念复查表}
\markboth{线性代数概念复查表}{线性代数概念复查表}

\begingroup
\setlength{\tabcolsep}{5pt}
\renewcommand{\arraystretch}{1.25}
\begin{longtable}{@{}>{\raggedright\arraybackslash}p{0.31\linewidth}>{\raggedright\arraybackslash}p{0.31\linewidth}>{\raggedright\arraybackslash}p{0.31\linewidth}@{}}
\hline
\endfirsthead
\hline
\endhead
\hline
\endfoot
\hline
\endlastfoot
\hyperref[p1:concept:c1-l7]{半群} & \hyperref[p1:concept:c7-l1976]{半正定二次型} & \hyperref[p1:concept:c7-l1976]{半正定矩阵}\\
\hyperref[p1:concept:c6-l2084]{伴随变换（伴随算子）} & \hyperref[p1:concept:c3-l2481]{伴随矩阵} & \hyperref[p1:concept:c1-l1154]{本原多项式}\\
\hyperref[p1:concept:c6-l46]{标准内积} & \hyperref[p1:concept:c2-l894]{标准形矩阵} & \hyperref[p1:concept:c6-l399]{标准正交基}\\
\hyperref[p1:concept:c4-l1313]{补子空间} & \hyperref[p1:concept:c5-l2373]{不变因子} & \hyperref[p1:concept:c5-l8]{不变子空间}\\
\hyperref[p1:concept:c7-l1972]{不定二次型} & \hyperref[p1:concept:c7-l1972]{不定矩阵} & \hyperref[p1:concept:c1-l558]{不定元}\\
\hyperref[p1:concept:c1-l789]{不可约多项式} & \hyperref[p1:concept:c1-l991]{超越数} & \hyperref[p1:concept:c7-l795]{初等合同变换}\\
\hyperref[p1:concept:c3-l920]{初等矩阵} & \hyperref[p1:concept:c5-l2454]{初等因子} & \hyperref[p1:concept:c5-l2058]{初等$\lambda$-矩阵}\\
\hyperref[p1:concept:c5-l2058]{初等$\lambda$-列变换} & \hyperref[p1:concept:c5-l2058]{初等$\lambda$-行变换} & \hyperref[p1:concept:c1-l1081]{代数闭域}\\
\hyperref[p1:concept:c1-l7]{代数结构} & \hyperref[p1:concept:c1-l991]{代数数} & \hyperref[p1:concept:c1-l1278]{代数无关}\\
\hyperref[p1:concept:c3-l2331]{代数余子式} & \hyperref[p1:concept:c1-l1024]{单根} & \hyperref[p1:concept:c6-l101]{单位化}\\
\hyperref[p1:concept:c3-l717]{单位矩阵} & \hyperref[p1:concept:c6-l101]{单位向量} & \hyperref[p1:concept:c1-l1278]{单项式}\\
\hyperref[p1:concept:c1-l1278]{单项式的次数} & \hyperref[p1:concept:c1-l1278]{单项式的系数} & \hyperref[p1:concept:c2-l1057]{导出方程组}\\
\hyperref[p1:concept:c2-l7]{等价的线性方程组} & \hyperref[p1:concept:c7-l1182]{等价二次型} & \hyperref[p1:concept:c2-l528]{等价向量组}\\
\hyperref[p1:concept:c6-l531]{等距的内积空间} & \hyperref[p1:concept:c6-l516]{等距映射} & \hyperref[p1:concept:c6-l493]{度量映射}\\
\hyperref[p1:concept:c6-l1545]{对称变换} & \hyperref[p1:concept:c1-l1460]{对称多项式} & \hyperref[p1:concept:c3-l856]{对称矩阵}\\
\hyperref[p1:concept:c7-l668]{对称双线性型} & \hyperref[p1:concept:c3-l1870]{对换} & \hyperref[p1:concept:c3-l908]{对角矩阵}\\
\hyperref[p1:concept:c7-l150]{对偶基} & \hyperref[p1:concept:c7-l12]{对偶空间} & \hyperref[p1:concept:c1-l926]{多项式的导数}\\
\hyperref[p1:concept:c1-l926]{多项式的高阶导数} & \hyperref[p1:concept:c1-l991]{多项式的根} & \hyperref[p1:concept:c1-l644]{多项式的公因子}\\
\hyperref[p1:concept:c1-l2482]{多项式的判别式} & \hyperref[p1:concept:c1-l789]{多项式的平凡因子} & \hyperref[p1:concept:c1-l558]{多项式的系数}\\
\hyperref[p1:concept:c1-l634]{多项式的因子} & \hyperref[p1:concept:c1-l634]{多项式的整除} & \hyperref[p1:concept:c1-l644]{多项式互素}\\
\hyperref[p1:concept:c1-l1278]{多元多项式} & \hyperref[p1:concept:c1-l1334]{多元多项式的乘法} & \hyperref[p1:concept:c1-l1315]{多元多项式的次数}\\
\hyperref[p1:concept:c1-l1334]{多元多项式的加法} & \hyperref[p1:concept:c1-l1324]{多元多项式的首项} & \hyperref[p1:concept:c1-l1428]{多元多项式函数}\\
\hyperref[p1:concept:c1-l1278]{多元多项式环} & \hyperref[p1:concept:c7-l1012]{二次型} & \hyperref[p1:concept:c7-l1031]{二次型的标准型}\\
\hyperref[p1:concept:c7-l1056]{二次型的矩阵} & \hyperref[p1:concept:c7-l1027]{二次型的系数} & \hyperref[p1:concept:c7-l1056]{二次型的秩}\\
\hyperref[p1:concept:c3-l856]{反对称矩阵} & \hyperref[p1:concept:c6-l809]{反射变换} & \hyperref[p1:concept:c2-l47]{方程组的初等行变换}\\
\hyperref[p1:concept:c2-l654]{方阵} & \hyperref[p1:concept:c5-l2373]{方阵的行列式因子} & \hyperref[p1:concept:c7-l392]{非退化双线性型}\\
\hyperref[p1:concept:c7-l1164]{非退化线性替换} & \hyperref[p1:concept:c1-l2000]{费马素数} & \hyperref[p1:concept:c3-l1322]{分块矩阵}\\
\hyperref[p1:concept:c3-l1440]{分块矩阵的乘法} & \hyperref[p1:concept:c3-l1440]{分块矩阵的加法} & \hyperref[p1:concept:c3-l1440]{分块矩阵的数乘}\\
\hyperref[p1:concept:c3-l1440]{分块矩阵的转置} & \hyperref[p1:concept:c7-l1527]{符号差} & \hyperref[p1:concept:c7-l1931]{负定二次型}\\
\hyperref[p1:concept:c7-l1931]{负定矩阵} & \hyperref[p1:concept:c7-l1527]{负惯性指数} & \hyperref[p1:concept:c7-l1425]{复二次型的规范型}\\
\hyperref[p1:concept:c3-l683]{复共轭转置} & \hyperref[p1:concept:c6-l2464]{复内积} & \hyperref[p1:concept:c1-l1024]{根的重数}\\
\hyperref[p1:concept:c1-l2164]{公倍式} & \hyperref[p1:concept:c1-l2305]{共轭多项式} & \hyperref[p1:concept:c4-l519]{过渡矩阵}\\
\hyperref[p1:concept:c6-l669]{合同（相合）矩阵} & \hyperref[p1:concept:c3-l363]{恒等映射} & \hyperref[p1:concept:c1-l7]{环}\\
\hyperref[p1:concept:c4-l227]{基} & \hyperref[p1:concept:c1-l1493]{基本对称多项式} & \hyperref[p1:concept:c3-l660]{基本矩阵}\\
\hyperref[p1:concept:c2-l1151]{基础解系} & \hyperref[p1:concept:c4-l227]{基向量} & \hyperref[p1:concept:c2-l548]{极大无关子组}\\
\hyperref[p1:concept:c5-l171]{几何重数} & \hyperref[p1:concept:c7-l1027]{交叉项} & \hyperref[p1:concept:c1-l7]{交换环}\\
\hyperref[p1:concept:c2-l61]{阶梯形方程组} & \hyperref[p1:concept:c2-l747]{阶梯形矩阵} & \hyperref[p1:concept:c3-l3286]{结式矩阵}\\
\hyperref[p1:concept:c2-l654]{矩阵} & \hyperref[p1:concept:c3-l504]{矩阵乘法} & \hyperref[p1:concept:c2-l857]{矩阵的初等变换}\\
\hyperref[p1:concept:c2-l857]{矩阵的初等列变换} & \hyperref[p1:concept:c2-l857]{矩阵的初等行变换} & \hyperref[p1:concept:c3-l826]{矩阵的迹}\\
\hyperref[p1:concept:c3-l787]{矩阵的极小多项式} & \hyperref[p1:concept:c3-l1322]{矩阵的块} & \hyperref[p1:concept:c3-l717]{矩阵的幂}\\
\hyperref[p1:concept:c3-l1987]{矩阵的特征多项式} & \hyperref[p1:concept:c2-l654]{矩阵的元素} & \hyperref[p1:concept:c2-l904]{矩阵的秩}\\
\hyperref[p1:concept:c3-l683]{矩阵的转置} & \hyperref[p1:concept:c3-l395]{矩阵加法} & \hyperref[p1:concept:c2-l654]{矩阵空间}\\
\hyperref[p1:concept:c3-l483]{矩阵数乘} & \hyperref[p1:concept:c5-l241]{可对角化的线性变换} & \hyperref[p1:concept:c5-l1191]{可对角化矩阵}\\
\hyperref[p1:concept:c3-l1016]{可逆矩阵} & \hyperref[p1:concept:c5-l1937]{可逆$\lambda$-矩阵} & \hyperref[p1:concept:c1-l789]{可约多项式}\\
\hyperref[p1:concept:c1-l526]{扩域} & \hyperref[p1:concept:c3-l246]{列满秩矩阵} & \hyperref[p1:concept:c2-l222]{列向量}\\
\hyperref[p1:concept:c2-l222]{列向量空间} & \hyperref[p1:concept:c2-l739]{列秩} & \hyperref[p1:concept:c4-l1459]{零度}\\
\hyperref[p1:concept:c1-l558]{零多项式} & \hyperref[p1:concept:c3-l757]{零化多项式} & \hyperref[p1:concept:c4-l75]{零空间}\\
\hyperref[p1:concept:c2-l222]{零向量} & \hyperref[p1:concept:c1-l7]{零因子} & \hyperref[p1:concept:c3-l7]{零映射}\\
\hyperref[p1:concept:c1-l1974]{梅森素数} & \hyperref[p1:concept:c3-l717]{幂零矩阵} & \hyperref[p1:concept:c6-l6]{内积}\\
\hyperref[p1:concept:c6-l572]{内积的度量矩阵} & \hyperref[p1:concept:c6-l6]{内积的对称性} & \hyperref[p1:concept:c6-l6]{内积的正定性}\\
\hyperref[p1:concept:c3-l1845]{逆对} & \hyperref[p1:concept:c3-l1016]{逆矩阵} & \hyperref[p1:concept:c3-l1845]{逆序数}\\
\hyperref[p1:concept:c4-l1979]{逆映射} & \hyperref[p1:concept:c1-l333]{辗转除法} & \hyperref[p1:concept:c1-l197]{欧几里德环}\\
\hyperref[p1:concept:c6-l6]{欧氏空间} & \hyperref[p1:concept:c3-l1845]{偶排列} & \hyperref[p1:concept:c3-l1812]{排列}\\
\hyperref[p1:concept:c3-l1845]{排列的符号} & \hyperref[p1:concept:c5-l22]{平凡不变子空间} & \hyperref[p1:concept:c4-l639]{平凡子空间}\\
\hyperref[p1:concept:c7-l1027]{平方项} & \hyperref[p1:concept:c1-l1379]{齐次多项式} & \hyperref[p1:concept:c2-l7]{齐次线性方程组}\\
\hyperref[p1:concept:c3-l1845]{奇排列} & \hyperref[p1:concept:c1-l7]{群} & \hyperref[p1:concept:c4-l276]{商空间}\\
\hyperref[p1:concept:c3-l881]{上三角矩阵} & \hyperref[p1:concept:c7-l1565]{实对称矩阵的负惯性指数} & \hyperref[p1:concept:c7-l1565]{实对称矩阵的正惯性指数}\\
\hyperref[p1:concept:c7-l1451]{实二次型的规范型} & \hyperref[p1:concept:c6-l6]{实内积空间} & \hyperref[p1:concept:c1-l558]{首系数}\\
\hyperref[p1:concept:c1-l558]{首一多项式} & \hyperref[p1:concept:c1-l518]{数域} & \hyperref[p1:concept:c6-l6]{双线性}\\
\hyperref[p1:concept:c7-l335]{双线性函数（双线性型）} & \hyperref[p1:concept:c7-l411]{双线性型的度量矩阵} & \hyperref[p1:concept:c7-l392]{双线性型的右根}\\
\hyperref[p1:concept:c7-l485]{双线性型的秩} & \hyperref[p1:concept:c7-l392]{双线性型的左根} & \hyperref[p1:concept:c7-l335]{双线性型空间}\\
\hyperref[p1:concept:c7-l2876]{双线性型下的正交补} & \hyperref[p1:concept:c7-l250]{双重对偶空间} & \hyperref[p1:concept:c7-l1848]{顺序主子式}\\
\hyperref[p1:concept:c1-l98]{素数} & \hyperref[p1:concept:c1-l498]{素因子的阶数} & \hyperref[p1:concept:c3-l2291]{特殊线性群}\\
\hyperref[p1:concept:c6-l1156]{特殊正交群} & \hyperref[p1:concept:c5-l1709]{特征矩阵} & \hyperref[p1:concept:c5-l171]{特征向量}\\
\hyperref[p1:concept:c5-l171]{特征值} & \hyperref[p1:concept:c5-l495]{特征值的代数重数} & \hyperref[p1:concept:c5-l171]{特征子空间}\\
\hyperref[p1:concept:c4-l2065]{同构的线性空间} & \hyperref[p1:concept:c1-l1278]{同类项} & \hyperref[p1:concept:c4-l3141]{凸集}\\
\hyperref[p1:concept:c4-l375]{维数} & \hyperref[p1:concept:c2-l7]{未知量} & \hyperref[p1:concept:c4-l375]{无穷维线性空间}\\
\hyperref[p1:concept:c2-l705]{系数矩阵} & \hyperref[p1:concept:c3-l881]{下三角矩阵} & \hyperref[p1:concept:c4-l3135]{线段}\\
\hyperref[p1:concept:c4-l1349]{线性变换} & \hyperref[p1:concept:c3-l363]{线性变换的多项式} & \hyperref[p1:concept:c3-l363]{线性变换的幂}\\
\hyperref[p1:concept:c5-l8]{线性变换的限制} & \hyperref[p1:concept:c4-l2846]{线性变换的最小多项式} & \hyperref[p1:concept:c2-l7]{线性方程组}\\
\hyperref[p1:concept:c2-l7]{线性方程组的解} & \hyperref[p1:concept:c7-l12]{线性函数} & \hyperref[p1:concept:c4-l7]{线性空间}\\
\hyperref[p1:concept:c7-l2637]{线性空间的张量积} & \hyperref[p1:concept:c4-l1981]{线性空间的自同构} & \hyperref[p1:concept:c4-l1981]{线性空间的自同构群}\\
\hyperref[p1:concept:c4-l722]{线性生成子空间} & \hyperref[p1:concept:c4-l1940]{线性同构映射} & \hyperref[p1:concept:c2-l270]{线性无关}\\
\hyperref[p1:concept:c4-l173]{线性无关集} & \hyperref[p1:concept:c2-l270]{线性相关} & \hyperref[p1:concept:c4-l173]{线性相关集}\\
\hyperref[p1:concept:c4-l1349]{线性映射} & \hyperref[p1:concept:c3-l328]{线性映射的复合} & \hyperref[p1:concept:c3-l209]{线性映射的核}\\
\hyperref[p1:concept:c3-l300]{线性映射的加法} & \hyperref[p1:concept:c3-l77]{线性映射的矩阵（标准基下）} & \hyperref[p1:concept:c3-l300]{线性映射的数乘}\\
\hyperref[p1:concept:c3-l248]{线性映射的像} & \hyperref[p1:concept:c7-l2824]{线性映射的张量积} & \hyperref[p1:concept:c4-l2617]{线性映射的秩}\\
\hyperref[p1:concept:c4-l1397]{线性映射下子集的像} & \hyperref[p1:concept:c4-l1391]{线性映射下子集的原像} & \hyperref[p1:concept:c4-l2360]{线性映射在给定基下的矩阵}\\
\hyperref[p1:concept:c4-l104]{线性组合} & \hyperref[p1:concept:c1-l558]{相伴多项式} & \hyperref[p1:concept:c5-l2001]{相抵的$\lambda$-矩阵}\\
\hyperref[p1:concept:c5-l820]{相似矩阵} & \hyperref[p1:concept:c2-l222]{向量} & \hyperref[p1:concept:c6-l3040]{向量到子空间的距离}\\
\hyperref[p1:concept:c4-l7]{向量的负元} & \hyperref[p1:concept:c2-l270]{向量的线性表示} & \hyperref[p1:concept:c7-l2632]{向量的张量积}\\
\hyperref[p1:concept:c4-l7]{向量加法} & \hyperref[p1:concept:c6-l157]{向量间的夹角} & \hyperref[p1:concept:c6-l101]{向量间的距离}\\
\hyperref[p1:concept:c4-l7]{向量数乘} & \hyperref[p1:concept:c6-l3075]{向量与子空间的夹角} & \hyperref[p1:concept:c6-l101]{向量正交}\\
\hyperref[p1:concept:c2-l383]{向量组的伸长组} & \hyperref[p1:concept:c2-l528]{向量组的线性表出} & \hyperref[p1:concept:c2-l548]{向量组的秩}\\
\hyperref[p1:concept:c6-l386]{向子空间的正交投影} & \hyperref[p1:concept:c2-l206]{行简化阶梯方程组} & \hyperref[p1:concept:c2-l873]{行简化阶梯形矩阵}\\
\hyperref[p1:concept:c3-l1908]{行列式} & \hyperref[p1:concept:c3-l263]{行满秩矩阵} & \hyperref[p1:concept:c2-l222]{行向量}\\
\hyperref[p1:concept:c2-l222]{行向量空间} & \hyperref[p1:concept:c2-l739]{行秩} & \hyperref[p1:concept:c5-l3660]{循环矩阵}\\
\hyperref[p1:concept:c5-l59]{循环子空间（Krylov空间）} & \hyperref[p1:concept:c2-l61]{严格阶梯形方程组} & \hyperref[p1:concept:c6-l245]{沿向量的正交投影}\\
\hyperref[p1:concept:c1-l7]{幺半群} & \hyperref[p1:concept:c3-l1171]{一般线性群} & \hyperref[p1:concept:c1-l558]{一元多项式}\\
\hyperref[p1:concept:c1-l573]{一元多项式的乘法} & \hyperref[p1:concept:c1-l558]{一元多项式的次数} & \hyperref[p1:concept:c1-l573]{一元多项式的加法}\\
\hyperref[p1:concept:c1-l573]{一元多项式的减法} & \hyperref[p1:concept:c1-l558]{一元多项式的首项} & \hyperref[p1:concept:c1-l991]{一元多项式函数}\\
\hyperref[p1:concept:c1-l558]{一元多项式环} & \hyperref[p1:concept:c4-l1707]{映射空间} & \hyperref[p1:concept:c4-l177]{由向量集合表出}\\
\hyperref[p1:concept:c7-l583]{由映射诱导的双线性型} & \hyperref[p1:concept:c5-l2716]{友矩阵} & \hyperref[p1:concept:c5-l2799]{有理标准型}\\
\hyperref[p1:concept:c5-l2759]{有理块} & \hyperref[p1:concept:c4-l375]{有限维线性空间} & \hyperref[p1:concept:c6-l2464]{酉空间}\\
\hyperref[p1:concept:c4-l708]{余维数} & \hyperref[p1:concept:c3-l2331]{余子式} & \hyperref[p1:concept:c1-l7]{域}\\
\hyperref[p1:concept:c2-l705]{增广矩阵} & \hyperref[p1:concept:c6-l101]{长度（范数）} & \hyperref[p1:concept:c1-l7]{整环}\\
\hyperref[p1:concept:c1-l101]{整数的公因子} & \hyperref[p1:concept:c1-l59]{整数的平凡因子} & \hyperref[p1:concept:c1-l59]{整数的因子}\\
\hyperref[p1:concept:c1-l59]{整数的整除} & \hyperref[p1:concept:c1-l101]{整数互素} & \hyperref[p1:concept:c7-l1578]{正定二次型}\\
\hyperref[p1:concept:c6-l698]{正定矩阵} & \hyperref[p1:concept:c7-l1623]{正定矩阵（二次型定义）} & \hyperref[p1:concept:c7-l1527]{正惯性指数}\\
\hyperref[p1:concept:c6-l2272]{正规变换（正规算子）} & \hyperref[p1:concept:c6-l2295]{正规矩阵} & \hyperref[p1:concept:c6-l789]{正交变换}\\
\hyperref[p1:concept:c6-l257]{正交补} & \hyperref[p1:concept:c6-l386]{正交分量} & \hyperref[p1:concept:c6-l981]{正交矩阵}\\
\hyperref[p1:concept:c6-l1043]{正交群} & \hyperref[p1:concept:c6-l1225]{正交相似} & \hyperref[p1:concept:c4-l1323]{直和的投影}\\
\hyperref[p1:concept:c1-l1278]{指数组} & \hyperref[p1:concept:c3-l1812]{置换群} & \hyperref[p1:concept:c1-l1024]{重根}\\
\hyperref[p1:concept:c1-l916]{重因子} & \hyperref[p1:concept:c1-l916]{重因子的重数} & \hyperref[p1:concept:c2-l61]{主元}\\
\hyperref[p1:concept:c3-l1322]{准对角矩阵} & \hyperref[p1:concept:c3-l1322]{准上三角矩阵} & \hyperref[p1:concept:c1-l1324]{字典序}\\
\hyperref[p1:concept:c3-l1812]{自然排列} & \hyperref[p1:concept:c3-l359]{自同态环} & \hyperref[p1:concept:c2-l129]{自由变量}\\
\hyperref[p1:concept:c3-l2544]{子矩阵} & \hyperref[p1:concept:c4-l585]{子空间} & \hyperref[p1:concept:c4-l881]{子空间的和}\\
\hyperref[p1:concept:c4-l276]{子空间的陪集} & \hyperref[p1:concept:c4-l1323]{子空间的嵌入} & \hyperref[p1:concept:c4-l722]{子空间的生成元}\\
\hyperref[p1:concept:c4-l1080]{子空间的直和} & \hyperref[p1:concept:c3-l2632]{子式} & \hyperref[p1:concept:c1-l526]{子域}\\
\hyperref[p1:concept:c1-l644]{最大公因式} & \hyperref[p1:concept:c1-l101]{最大公因子} & \hyperref[p1:concept:c1-l2164]{最小公倍式}\\
\hyperref[p1:concept:c4-l502]{坐标} & \hyperref[p1:concept:c4-l502]{坐标分量} & \hyperref[p1:concept:c7-l108]{坐标函数}\\
\hyperref[p1:concept:c3-l7]{坐标向量空间上的线性映射} & \hyperref[p1:concept:c3-l7]{坐标向量空间上的End空间} & \hyperref[p1:concept:c3-l7]{坐标向量空间上的Hom空间}\\
\hyperref[p1:concept:c4-l502]{坐标映射} & \hyperref[p1:concept:c1-l7]{Abel群} & \hyperref[p1:concept:c4-l1349]{End空间}\\
\hyperref[p1:concept:c6-l624]{Gram矩阵} & \hyperref[p1:concept:c4-l3090]{Hamel基} & \hyperref[p1:concept:c4-l1349]{Hom空间}\\
\hyperref[p1:concept:c5-l2950]{Jordan标准型} & \hyperref[p1:concept:c5-l2911]{Jordan矩阵} & \hyperref[p1:concept:c5-l2911]{Jordan块}\\
\hyperref[p1:concept:c6-l399]{Kronecker符号} & \hyperref[p1:concept:c3-l2915]{Kronecker积（张量积）} & \hyperref[p1:concept:c3-l652]{Kronecker记号}\\
\hyperref[p1:concept:c5-l1709]{$\lambda$-矩阵} & \hyperref[p1:concept:c5-l2058]{$\lambda$-矩阵的初等变换} & \hyperref[p1:concept:c5-l1937]{$\lambda$-矩阵的逆}\\
\hyperref[p1:concept:c5-l1819]{$\lambda$-矩阵的行列式} & \hyperref[p1:concept:c5-l2205]{$\lambda$-矩阵的行列式因子} & \hyperref[p1:concept:c5-l1905]{$\lambda$-矩阵的秩}\\
\hyperref[p1:concept:c1-l1385]{Newton和} & \hyperref[p1:concept:c6-l2744]{QR分解} & \hyperref[p1:concept:c4-l3113]{Schauder基}\\
\hyperref[p1:concept:c5-l2095]{Smith标准形} &  & \\
\end{longtable}
\endgroup
